\section{Actores e interacción}
\label{sec:interaccion}

\subsection{Roles}

\begin{table}[H]
\centering\small
\begin{tabularx}{\textwidth}{@{}l L L@{}}
\toprule
\textbf{Rol} & \textbf{Qué necesita saber} & \textbf{Cuándo interviene} \\
\midrule
Guardia de la central & Dónde está el robot, qué puntos verificó, si hay alertas abiertas y si los datos están al día & Inicia la ronda, decide ante un incidente, reconoce alertas, manda el robot a la base \\
Jefe de seguridad & Qué rondas se hicieron, cuándo y qué pasó en cada una & Al día siguiente, revisando el historial \\
Personas en el edificio & Qué es el robot y qué tienen que hacer cuando lo cruzan & Cuando se lo encuentran en un pasillo o en una zona restringida \\
\bottomrule
\end{tabularx}
\end{table}

\subsection{HMI: lo que ve el guardia}

La pantalla principal es la de supervisión. Arriba está el estado de la conexión con la hora del último dato; en el centro, el plano con la ruta, los puntos ya verificados y la posición del robot; abajo, las órdenes disponibles. Cada orden muestra si se puede usar y, si no, por qué.

Usamos tema oscuro por defecto. El guardia trabaja de noche en una central con poca luz y una pantalla blanca encandila. Los colores tienen un solo significado en toda la app: verde para lo normal, ámbar para advertencias, rojo para errores y gris para datos viejos.

\subsection{HRI: lo que ven las personas que se cruzan con el robot}

El robot no está pensado para hablar con la gente, pero se la va a cruzar: personal de limpieza, un docente que se quedó tarde, alguien de mantenimiento. Para esos casos definimos un comportamiento fijo:

\begin{itemize}
  \item Si detecta a una persona en un lugar de paso, se detiene a dos metros, le da paso y sigue.
  \item Si la persona está en una zona restringida, se detiene a distancia y dice por el parlante: ``Zona restringida. Seguridad fue notificada''. No la sigue ni la encierra.
  \item El robot nunca toma decisiones sobre la persona. Avisa y es el guardia quien la clasifica.
\end{itemize}

\subsection{Cuándo tiene que intervenir un humano}

El robot sigue la ronda solo, pero frena y espera al guardia en cuatro situaciones: algo le bloquea el camino, detecta a una persona en una zona restringida, la batería baja del umbral crítico o pierde la conexión por más de un minuto. En este último caso no puede esperar a nadie, así que vuelve solo a la base.
